\section{十三、内存属性与 A/D 位：翻译结果不只是页框号}

地址转换常被简写为“VPN 换成 PFN”，但处理器真正得到的是一组访问属性：物理页框、读写执行权限、用户或特权属性、页大小、全局或地址空间身份，以及可缓存性和排序类型。后续 Cache、互连与设备桥必须使用同一组结果。若只让页框号正确而属性错误，同一物理地址仍可能表现为两份不一致的 Cache 数据，或把一次设备寄存器访问变成推测、合并和重复访问。

\begin{pdfdiagram}[x=1cm,y=1cm]
  \node[hw state, minimum width=2.3cm] (va) at (1.3,3.8) {虚拟地址访问\\load/store/fetch};
  \node[hw control, minimum width=2.6cm] (walk) at (4.2,3.8) {TLB 或 page walk\\累计各级约束};
  \node[hw memory, minimum width=3.0cm] (tuple) at (7.6,3.8) {翻译结果\\PFN、权限、页大小\\内存类型、ASID};
  \node[hw control, minimum width=2.7cm] (check) at (11.1,3.8) {权限与类型检查\\决定后续路径};
  \draw[hw bus] (va) -- (walk);
  \draw[hw bus] (walk) -- (tuple);
  \draw[hw bus] (tuple) -- (check);

  \node[hw compute, minimum width=3.0cm] (normal) at (8.8,1.0) {普通内存路径\\Cache、预取、一致性};
  \node[hw block, minimum width=3.0cm] (device) at (12.3,1.0) {设备内存路径\\受限宽度、顺序、副作用};
  \coordinate (typebranch) at (11.1,2.25);
  \draw[BookMuted, line width=1.0pt] (check.south) -- (typebranch);
  \draw[hw bus] (typebranch) -- (normal.north);
  \draw[hw bus] (typebranch) -- (12.3,2.25) -- (device.north);
\end{pdfdiagram}
\diagramnote{一次地址翻译产生的是“页框加访问属性”元组。普通内存可以进入 Cache 与推测路径，设备内存必须遵守更严格的宽度、顺序和副作用规则；两条路径在翻译结果处已经分开。}

\topic{Accessed 与 Dirty 位是硬件和操作系统共享的状态}

Accessed 位记录页表层级或叶子映射曾参与访问，Dirty 位记录叶子页曾被写。操作系统可以周期性清除这些位，过一段时间再观察哪些页面重新置位，从而估计工作集、选择回收候选并判断匿名页或文件页是否需要回写。位为 0 只表示“自最近一次清除以来尚未观察到对应事件”，不能证明页面历史上从未访问。

硬件置位与软件清位可能并发。假设操作系统读到 A=1，准备清零；此时另一核心又访问该页。如果软件用普通读—改—写覆盖整个 PTE，可能把硬件刚写入的 A 或 D 位丢失。体系结构因此规定页表项更新的原子粒度和失效要求，内核使用原子操作、页表锁或再次检查来闭合竞争。若某体系结构不由硬件置 A/D，操作系统可先把映射设成会触发访问异常的状态，在首次读或写时由异常处理程序记录并放开权限；代价是第一次访问增加一次受控 fault。

清除 A/D 位不一定改变页框号，却可能仍需要 TLB 处理。TLB 中若缓存了旧访问状态或允许硬件只在首次填入时更新页表，软件必须按体系结构规定失效对应翻译，才能保证后续访问重新反映到内存中的 PTE。只改页表内存而不处理已缓存翻译，会得到长期保持为 0 的假“冷页”。

\topic{内存类型决定推测、合并与可见顺序}

普通回写内存允许 Cache line 填充、预取、写合并和一定程度的乱序；写穿或不可缓存内存使用不同的带宽与可见性规则；设备内存还可能禁止推测读取、限制访问宽度并要求程序顺序。设备状态寄存器的读取可能具有“读后清零”副作用，doorbell 写入会触发外部动作。若把这类页面映射成普通可缓存内存，CPU 可能减少、延后或合并总线事务，使寄存器语义失真。

同一物理页通过两个虚拟地址映射时，内存类型必须兼容。一个别名按回写 Cache 访问，另一个别名按不可缓存访问，可能让一个核心读取 Cache 中的旧副本，而设备或另一个别名已经修改物理内存。修复不能只执行一次普通屏障；系统要先停止相关访问，按平台规定清理和失效 Cache，统一属性后再重新建立映射。

\begin{longtable}{L{0.18\linewidth}L{0.28\linewidth}L{0.28\linewidth}L{0.16\linewidth}}
\toprule
属性维度 & 决定的问题 & 错误设置的典型后果 & 主要证据 \\
\midrule
R/W/X 与特权 & 指令或数据访问是否允许 & protection fault 或意外可执行页面 & fault 地址与错误码 \\\tablerowrule
A/D 位 & 页面近期是否读写、是否需回写 & 热页被误回收或脏页遗漏 & PTE 位与回收 trace \\\tablerowrule
普通/设备内存 & 能否 Cache、推测、合并和重排 & MMIO 副作用次数或顺序错误 & 页属性与总线 trace \\\tablerowrule
页大小 & TLB 覆盖、碎片与迁移粒度 & TLB miss 或大页拆分抖动 & TLB 事件与页统计 \\\tablerowrule
ASID/全局属性 & 翻译在哪些上下文中可复用 & 切换后误命中旧地址空间 & ASID generation 与失效记录 \\
\bottomrule
\end{longtable}

\section{十四、大页拆合与 NUMA 页面迁移}

大页并非创建后永久不变。操作系统可以把一段连续、权限一致且访问稳定的 4 KiB 页面合并为大页，扩大 TLB reach；也可能因为局部 COW、细粒度保护、内存压力或迁移要求把大页拆回小页。合并前要确认物理页连续、映射属性相同、没有不兼容的固定页或设备映射；发布大页表项后，还要失效旧的小页 TLB 项。拆分时则先准备完整的小页表，再原子替换上层大页表项并执行对应范围失效，不能让 walker 看到半初始化页表。

\topic{页面迁移转移的是内容、映射和页框生命周期}

NUMA 系统可以把频繁由某个处理器访问的页面迁到更近的内存节点。迁移并不是简单复制：内核先阻止映射产生新的可写访问或把 PTE 改成迁移状态，等待已有页表使用者进入可控边界；随后分配目标页框并复制数据，原子发布指向新页框的 PTE，向相关核心执行 TLB shootdown，最后才降低旧页引用计数并释放旧页框。

文件页、匿名页、共享页和固定页的迁移条件不同。文件页可从后备文件重新读取，但脏页要保持回写状态；匿名页没有文件副本，复制失败必须保留旧页；多个进程共享的页面要更新全部反向映射；正在进行 DMA 的固定页不能在设备仍持有旧 IOVA 或 ATS 翻译时移动。只更新 CPU PTE 而遗漏 IOMMU 和设备翻译，会让 DMA 继续写旧页框。

\topic{迁移收益必须覆盖复制、失效与重新预热成本}

将 2 MiB 页面跨节点复制，若有效复制带宽为 20 GiB/s，纯数据搬运至少约为 $2\,\mathrm{MiB}/20\,\mathrm{GiB/s}\approx98\,\mu\mathrm{s}$，还未计页表锁、TLB shootdown、Cache 失效和调度延迟。若页面之后只被远端访问几次，迁移成本可能高于节省的访问时延；只有页面保持足够热、访问者位置相对稳定时，迁移才有收益。

自动 NUMA 平衡通常先用采样 fault 或访问计数估计“谁在访问”，再设置阈值与安静期，避免页面在两个节点之间反复移动。大页使每次迁移覆盖更多数据，也可能因局部热点触发不必要复制；系统可先拆页，再只移动真正活跃的小页。评价策略时应同时观察远端访问比例、迁移字节、fault 数、shootdown 时延和迁移后的驻留时间。

\begin{longtable}{L{0.17\linewidth}L{0.28\linewidth}L{0.31\linewidth}L{0.14\linewidth}}
\toprule
操作 & 发布前必须成立 & 发布后必须完成 & 失败处理 \\
\midrule
小页合并为大页 & 物理连续、属性一致、无不兼容固定页 & 失效小页翻译并验证大页命中 & 保留原小页映射 \\\tablerowrule
大页拆分 & 小页表全部初始化完成 & 替换上层项并失效大页翻译 & 继续使用原大页 \\\tablerowrule
NUMA 页面迁移 & 新写受控、目标页已复制完整 & 发布新 PTE、shootdown、释放旧页 & 回滚到旧页框 \\\tablerowrule
设备共享页迁移 & 停止提交并耗尽在途 DMA & IOTLB/ATS 失效确认后才复用旧页 & 固定旧页或设备复位 \\
\bottomrule
\end{longtable}
